\chapter{Quick command reference}
\label{app:comandos}
\lstset{aboveskip=0.35em,belowskip=0.4em}

This appendix collects the main Linux commands in the procedure. The corresponding
chapters explain their use and conditions.

\section{Check the system}

\begin{lstlisting}[language=bash]
uname -r
uname -m
cat /etc/os-release
lspci -nn
\end{lstlisting}

\section{Raspberry Pi 5}

For a non-HAT+ PCIe adapter, edit:

\begin{lstlisting}[language=bash]
sudoedit /boot/firmware/config.txt
\end{lstlisting}

and add this line if it is not already present:

\begin{lstlisting}
dtparam=pciex1
\end{lstlisting}

Then reboot and check the connection:

\begin{lstlisting}[language=bash]
sudo reboot
lspci -nn
\end{lstlisting}

\section{PCIe driver}

On Ubuntu:

\begin{lstlisting}[language=bash]
sudo apt update
sudo apt install dkms build-essential "linux-headers-$(uname -r)"
\end{lstlisting}

Download and install the driver:

\begin{lstlisting}[language=bash]
git clone https://github.com/Brainchip-Inc/akida_dw_edma.git
cd akida_dw_edma
git rev-parse HEAD
./install.sh
\end{lstlisting}

Check the installation:

\begin{lstlisting}[language=bash]
dkms status
lsmod | grep -i akida
ls -l /dev/akida*
lspci -nnk
\end{lstlisting}

\section{Python environment}

From the project root:

\begin{lstlisting}[language=bash]
python3.11 -m venv .venv
source .venv/bin/activate
python -m pip install --upgrade pip
\end{lstlisting}

Current MetaTF installation:

\begin{lstlisting}[language=bash]
python -m pip install "metatf==2.19.3"
\end{lstlisting}

or the repository's reference environment:

\begin{lstlisting}[language=bash]
python -m pip install -r requirements/reproducible.txt
\end{lstlisting}

Then install the project tools:

\begin{lstlisting}[language=bash]
python -m pip install -e . --no-deps
python -m pip check
\end{lstlisting}

\section{Smoke tests}

Software test:

\begin{lstlisting}[language=bash]
python examples/00_simulator_smoke.py
akd1000-bringup verify \
  --id smoke \
  --model models/akida_v1_smoke_test.fbz
\end{lstlisting}

Diagnostics:

\begin{lstlisting}[language=bash]
akd1000-bringup doctor
\end{lstlisting}

Board test:

\begin{lstlisting}[language=bash]
akida devices
python examples/01_hardware_smoke.py
\end{lstlisting}

\section{Run your own model}

\begin{lstlisting}[language=bash]
akd1000-bringup run \
  --model models/model.fbz \
  --input data/input.npy \
  --method forward \
  --clock Performance \
  --map AllNps \
  --output outputs/raw-output.npy
\end{lstlisting}

\section{Benchmark}

\begin{lstlisting}[language=bash]
akd1000-bringup benchmark \
  --model models/model.fbz \
  --input data/inputs.npy \
  --method forward \
  --clock Performance \
  --map AllNps \
  --warmup 8 \
  --repetitions 1 \
  --output outputs/smoke
\end{lstlisting}

\section{Quick troubleshooting}

\begin{enumerate}
  \item If the board does not appear in \code{lspci}, check the assembly and PCIe configuration.
  \item If it appears on PCIe but \code{/dev/akida*} is missing, check the driver.
  \item If the node exists but \code{akida.devices()} is empty, check the environment and permissions.
  \item If the test model works but your own model does not, check the model and its input data.
\end{enumerate}
